\chapter{Modelagem do Transporte de Vacinas}

Os modelos propostos não contemplam a etapa inicial de distribuição de vacinas a partir do nível federal. Assume-se que os estoques iniciais nos municípios já refletem a alocação realizada pelo Ministério da Saúde e o modelo começa depois disso. Dessa forma, o problema abordado concentra-se na redistribuição eficiente de vacinas entre municípios, considerando desequilíbrios locais de oferta e demanda.

\section{Modelo Intermunicipal Sem Intermediação Estadual}

Este modelo representa um cenário hipotético de operação descentralizada, na qual no qual municípios podem redistribuir vacinas diretamente entre si sem a necessidade de passagem por órgãos intermediários, como secretarias estaduais de saúde, conforme representado a seguir:

\[
\text{Município } i \;\longrightarrow\; \text{Município } j
\]

A adoção dessa abordagem tem como principal objetivo simplificar a estrutura do modelo e, ao mesmo tempo, possibilitar avaliar o desempenho, em termos de tempo, de uma estratégia de distribuição direta.

\subsection{Variável de Decisão e Parâmetros}

A \textbf{variável de decisão} é definida:
\\
\\
\begin{equation*}
    x_{ijv} \in \mathbb{Z}_{\geq 0} \quad \forall i \in I, j \in J, t \in V,
\end{equation*}
onde:
\begin{itemize}
    \item $I$ representa os municípios $i$ mapeados;
    \item $J$ representa os municípios $j$ mapeados;
    \item $v$ representa os tipos de vacina mapeados.
    \item $x$ representa a quantidade de vacinas transportadas;
    \item $i$ indica o município de origem $i$, que possui oferta da vacina que será transportada;
    \item $j$ indica o município de destino $j$, que possui demanda pela vacina que será transportada;
    \item $v$ indica o tipo de vacina que será transportada, ou seja, definido por uma tabela indicando o índice apropriado para cada tipo. Por exemplo:
\end{itemize}

\begin{quadro}[htb] 
\centering
\caption{Exemplo de uso da variável $v$}
\label{quad:quadro1}
% Trocamos tabularx por tabular e removemos a largura fixa
\begin{tabular}{|c|l|} 
\hline
\textbf{$v$} & \textbf{Nome da vacina} \\
\hline
1 & BCG \\
\hline
2 & HEPATITE B \\
\hline
\end{tabular}
\end{quadro}

Além disso, como x representa a quantidade transportada, foi definido que seu valor deve ser um número inteiro e não-negativo, pois não faria sentido transportar meia dose de vacina, por exemplo. Desta forma, se trata de um \textbf{Problema de Programação Linear Inteira}.

Além da variável de decisão, é importante definir alguns \textbf{parâmetros}.

As principais opções reais de modais para o transporte de vacinas são o rodoviário e o aéreo e por isso eles serão considerados na modelagem. Esses modais possuem tempos de deslocamento diferentes para uma mesma distância, por isso se faz necessário definir esses \textbf{parâmetros de tempo}.

Sejam:

\[
\tau_{ij}^{rod}, \quad \tau_{ij}^{areo} \in \mathbb{R}_{\geq 0}
\]

onde:
\begin{itemize}
    \item $\tau_{ij}^{rod}$ é o tempo de deslocamento do modal rodoviário entre o município de origem $i$ e o município de destino $j$.
    \item $\tau_{ij}^{areo}$ é o tempo de deslocamento do modal aéreo entre o município de origem $i$ e o município de destino $j$.
\end{itemize}

No caso do transporte aéreo, o modelo não considera explicitamente apenas o tempo de voo entre os municípios. O tempo $\tau_{ij}^{areo}$ representa o tempo total agregado da operação logística aérea, incluindo:

\begin{itemize}
    \item transporte rodoviário do município de origem até o aeroporto;
    \item tempo de voo entre aeroportos;
    \item transporte rodoviário do aeroporto de destino até o município final.
\end{itemize}

Dessa forma, $\tau_{ij}^{areo}$ já incorpora implicitamente todas as etapas necessárias para a entrega via modal aéreo, mantendo o modelo compacto sem perda significativa de realismo.

Na prática, o processo de distribuição de vacinas não se limita apenas ao tempo de deslocamento entre os pontos de origem e destino. Ao longo da cadeia logística, existem etapas intermediárias que incluem atividades como desembarque, descarga, conferência, armazenamento temporário e redistribuição das vacinas.

Essas atividades introduzem um tempo adicional no processo, denominado neste trabalho como \textbf{tempo de processamento logístico} e tal tempo ocorre, principalmente, nos pontos de destino, como aeroportos, secretarias municipais de saúde ou centros de distribuição locais.

O tempo de processamento logístico pode variar de acordo com o modal de transporte utilizado. No caso do transporte aéreo, por exemplo, há etapas adicionais como operações aeroportuárias, inspeções e manuseio especializado, que tendem a aumentar o tempo total de processamento. Já no transporte rodoviário, o fluxo tende a ser mais direto, resultando em tempos menores. Para capturar esse comportamento, são considerados tempos de processamento logístico distintos por modal.

Modelar explicitamente essas etapas exigiria a introdução de novas variáveis de decisão e restrições, caracterizando um problema mais complexo e poderia comprometer a tratabilidade computacional. Em vez disso, os tempos de deslocamento e de processamento logístico e a escolha do modal mais adequado (rodoviário ou aéreo) serão incorporados, implicitamente, a \textbf{parâmetros de custo total}.

Sejam:

\[
\delta^{rod}, \quad \delta^{areo} \in \mathbb{R}_{\geq 0}
\]

onde:
\begin{itemize}
    \item $\delta^{rod}$ representa o tempo de processamento logístico do modal rodoviário.
    \item $\delta^{areo}$ representa o tempo de processamento logístico do modal aéreo.
\end{itemize}

Ainda, é necessário considerar as diferenças de custo financeiro e de complexidade logística entre os dois tipos de modais. O transporte aéreo é mais rápido, porém mais caro financeiramente e mais complexo logisticamente.

Seja:

\[
\alpha > 1
\]

um \textbf{fator de penalização associado ao transporte aéreo} e que reflete limitações operacionais e custos adicionais do transporte aéreo. 

Essa abordagem permite modelar a preferência pelo transporte rodoviário em situações onde ambos os modais são viáveis, ao mesmo tempo em que mantém o transporte aéreo como alternativa para cenários em que o ganho de tempo é significativo. Essa abordagem permite capturar as diferenças operacionais entre modais sem aumentar a complexidade do modelo.

Sejam:
\begin{itemize}
    \item $C_{ij}^{rod}$, o \textbf{Custo Total do transporte rodoviário};
    \item $C_{ij}^{areo}$, o \textbf{Custo Total do transporte aéreo}.
\end{itemize}


Dessa forma, o \textbf{Custo Total de transporte por modal} é dado por:

\[
C_{ij}^{rod} = \tau_{ij}^{rod} + \delta^{rod}
\]
\[
C_{ij}^{areo} = \alpha \cdot \tau_{ij}^{areo} + \delta^{areo}
\]


Dessa forma, define-se o \textbf{Custo Total de transporte} como:

\[
C_{ij} = \min \left( \quad C_{ij}^{rod} \quad , \quad C_{ij}^{areo} \quad \right),  \quad C_{ij} \in \mathbb{R}_{\geq 0}
\]

Assim, o custo total de transporte $C_{ij}$ é definido como o menor tempo efetivo entre os modais disponíveis.

Essa abordagem não introduz novas variáveis de decisão e, dessa forma, mantém a eficiência computacional do modelo, ao mesmo tempo em que aumenta sua aderência à realidade operacional do sistema de distribuição de vacinas.

Além das diferenças de tempo entre os modais de transporte, é necessário considerar que a \textbf{capacidade de transporte} também pode variar de acordo com o modal utilizado.

Sejam:
\[
u_{ij}^{\text{rod}}, \quad u_{ij}^{\text{areo}} \in \mathbb{R}_{\geq 0}
\]

onde:
\begin{itemize}
    \item $u_{ij}^{\text{rod}}$ representa a capacidade máxima agregada de transporte entre os municípios $i$ e $j$ utilizando o modal rodoviário;
    \item $u_{ij}^{\text{areo}}$ representa a capacidade máxima agregada de transporte entre os municípios $i$ e $j$ utilizando o modal aéreo.
\end{itemize}

Neste trabalho, a capacidade $u_{ab}$ não descreve o limite de um único veículo isolado, mas sim a \textbf{capacidade agregada de fluxo logístico} viabilizada pelo modal ao longo do horizonte de planejamento da campanha de distribuição. Essa abstração engloba a capacidade de um veículo multiplicada pelo tamanho da frota alocada para a rota e pelo número de viagens (ciclos) realizados no período.

Dessa forma, a capacidade efetiva de transporte entre os municípios $i$ e $j$ é definida de forma consistente com o custo escolhido, conforme:

\begin{equation}
u_{ij} =
\begin{cases}
u_{ij}^{\text{rod}}, & \text{se } C_{ij}^{\text{rod}} \le C_{ij}^{\text{areo}} \\
u_{ij}^{\text{areo}}, & \text{caso contrário}
\end{cases}
\end{equation}

Essa modelagem mantém a simplicidade do problema, evitando o aumento do número de variáveis de decisão, ao mesmo tempo em que incorpora diferenças operacionais relevantes entre os modais de transporte.

Com a variável de decisão e parâmetros definidos, pode-se definir o tempo de execução do plano de distribuição de vacinas, ou seja, quanto tempo levou até a última entrega chegar desde que a distribuição foi iniciada.

Considerando que as entregas intermunicipais ocorrem em paralelo, o tempo total de distribuição é determinado pela entrega que leva mais tempo a ser concluída. Dessa forma, o \textbf{tempo de distribuição} corresponde ao maior tempo de transporte entre todos os fluxos realizados no sistema.

Assim, o tempo de distribuição é dado por:


\[
T = \max \left\{ C_{ij} \;\middle|\; x_{ijv} > 0 \right\}
\]

onde:
\begin{itemize}
    \item $T$ é o tempo de distribuição das vacinas.
\end{itemize}

\subsection{Função Objetivo}

Assim, a \textbf{função objetivo} ficou definida da seguinte forma:

\begin{equation}
    \min z = \sum_{v} \sum_{i} \sum_{j} C_{ij} \cdot x_{ijv}
    \label{eq:intermunicipal_sem_penalidade}
\end{equation}

Portanto, o problema modelado tem como objetivo \textbf{minimizar o tempo necessário para entregar a quantidade necessária das doses de vacinas ofertadas para os municípios com demanda}.

Pode-se, ainda, fazer uma modificação da função anterior adicionando uma \textbf{penalidade $\rho_{iv}$} para a vacina de tipo $v$ do município $i$ de acordo com o tempo de expiração, para que, no processo de transporte, haja o menor número de doses expiradas possível. Assim, temos:


\[ 0 < \rho_{iv} \le 1, \]  \\
onde:
\begin{itemize}
    \item dado um tempo de expiração $e_{iv}$, ou seja, tempo para a vacina $v$ retida no município $i$ expirar, então a penalidade é definida por: $\rho_{iv} = \frac{1}{e_{iv}}$.
\end{itemize}

Assim, a função objetivo modificada fica definida da seguinte forma:

\begin{equation}
    \min z = \sum_{v} \sum_{i} \sum_{j} C_{ij} \cdot x_{ijv} \cdot \rho_{iv}
    \label{eq:intermunicipal_com_penalidade}
\end{equation}

Esta função tem um objetivo ligeiramente diferente: \textbf{minimizar o tempo ponderado pelo risco de expiração}.

Finalmente, algumas restrições precisam ser levadas em consideração para este problema.

\subsection{Restrições}

\textbf{A primeira restrição} é a \textbf{restrição de oferta} de cada vacina $v$ em cada município $i$. 
O número de vacinas transportadas pode ser menor que número de vacinas ofertadas no local de origem e, no máximo, pode-se transportar uma quantidade igual a quantidade ofertada pela origem. Assim, temos:

\begin{equation*}
    \sum_{j} x_{ijv} \le o_{iv}, \quad \forall i \in I, \; v \in V
\end{equation*}
onde:
\begin{itemize}
    \item $o_{ik}$ representa a quantidade de vacinas do tipo $v$ ofertadas no município $i$;
\end{itemize}

Seria um erro se a restrição fosse definida desta forma:
\begin{equation*}
    \sum_{j} x_{ijv} = o_{iv}, \quad \forall i \in I, \; v \in V
\end{equation*}
pois, nesse caso, todas as vacinas ofertadas seriam obrigadas a serem transportadas e, num cenário real, do conjunto de vacinas ofertadas pode sobrar vacina pode ou não haver destino viável.

\textbf{A segunda restrição} é a \textbf{restrição de demanda} da vacina $v$ por cada município $j$. 
O número de vacinas transportadas deve ser exatamente a quantidade demandada para que não haja falta ou desperdício. Assim, temos:

\begin{equation*}
    \sum_{i} x_{ijv} = d_{jv}, \quad \forall j \in J, \; v \in V
\end{equation*}
onde:
\begin{itemize}
    \item $d_{jv}$ representa a quantidade de vacinas do tipo $v$ demandadas no município $j$;
    \item $J$ representa os municípios $j$ mapeados.
\end{itemize}

E, nesse caso, estamos forçando o atendimento total da demanda.

Seria um erro se a restrição fosse definida desta forma:
\begin{equation*}
    \sum_{i} x_{ijv} \le d_{jv}, \quad \forall j \in J, \; v \in V
\end{equation*}
pois isso implicaria que se pode não atender a demanda e o modelo pode simplesmente não entregar nada (minimiza custo = 0).

\textbf{A terceira restrição} é a \textbf{restrição de capacidade agregada} de entrega de vacinas da origem $i$ ao destino $j$. 
O número de vacinas transportadas não pode exceder a capacidade máxima agregada. Assim, temos:

\begin{equation*}
    \sum_{v} x_{ijv} \le u_{ij}, \quad \forall i \in I, j \in J
\end{equation*}
onde:
\begin{itemize}
    \item $u_{ij}$ representa a capacidade máxima (limite superior) de vacinas que pode ser transportada do município $i$ para o município $j$.
\end{itemize}

Assume-se que diferentes tipos de vacina compartilham a mesma capacidade de transporte, o que é consistente com a operação logística real, na qual múltiplos imunizantes são transportados conjuntamente no mesmo veículo.

\textbf{A quarta restrição} é a \textbf{restrição de viabilidade temporal} e é relativa à ideia de remover pares onde o tempo de transporte é maior que o tempo restante até o vencimento. Assim, temos:

\begin{equation*}
x_{ijv} = 0, \quad \forall i \in I, \; j \in J, \; v \in V \text{ tais que } C_{ij} > e_{iv}
\end{equation*}

Em suma, o PPLI ficou definido assim:

\begin{equation*}
    \min z = \sum_{v} \sum_{i} \sum_{j} C_{ij} \cdot x_{ijv}
\end{equation*}
sujeito a:
\begin{equation*}
    \sum_{j} x_{ijv} \le o_{ik}, \quad \forall i \in I, \; v \in V
\end{equation*}

\begin{equation*}
    \sum_{i} x_{ijv} = d_{jv}, \quad \forall j \in J, \; v \in V
\end{equation*}

\begin{equation*}
    \sum_{v} x_{ijv} \le u_{ij}, \quad \forall i \in I, j \in J
\end{equation*}

\begin{equation*}
x_{ijv} = 0, \quad \forall i \in I, \; j \in J, \; v \in V \text{ tais que } C_{ij} > e_{iv}
\end{equation*}

\begin{equation*}
    x_{ijv} \in \mathbb{Z}_{\geq 0}, \quad \forall i \in I, j \in J, v \in V \quad \text{(restrição de não-negatividade)}
\end{equation*}


\section{Modelo Intermunicipal Com Intermediação Estadual}

Este modelo busca representar o trajeto físico real considerando o intermédio de órgãos estaduais (como a secretaria estadual) no transporte de vacinas entre municípios.

Usaremos
\[ k \in K \]
para representar o centro estadual $k$ dentre os centros estaduais $K$ mapeados.

Com isso, o fluxo pode ser representado da seguinte forma:

\[
\text{Município } i \;\longrightarrow\; \text{Estado } k \;\longrightarrow\; \text{Município } j
\]

\subsection{Variável de Decisão e Parâmetros}

Para o transporte, são necessárias duas etapas de fluxo:

\begin{enumerate}
    \item Município $i \longrightarrow$ Estado $k$ \\
    Para essa etapa, temos:
    \[ x_{ikv} \quad \text{e} \quad C_{ik}\]
    onde:
    \begin{itemize}
    \item $x_{ikv}$ representa a quantidade de vacinas do tipo $v$ transportadas do município $i$ ao centro estadual $k$;
    \item $C_{ik}$ representa o \textbf{custo total} (em horas) para se transportar vacinas do município $i$ ao centro estadual $k$.
    \end{itemize}.
    
    \item Estado $k \longrightarrow$ Município $j$ \\
    Para essa etapa, temos:
    \[y_{kjv} \quad \text{e} \quad C_{kj}\]
    onde:
     \begin{itemize}
    \item $y_{kjv}$ representa a quantidade de vacinas do tipo $v$ transportadas do centro estadual $k$ ao município $j$.
    \item $C_{kj}$ representa o \textbf{custo total} (em horas) para se transportar vacinas do centro estadual $k$ ao município $j$.
    \end{itemize}.
\end{enumerate}

Foram definidas duas variáveis de decisão distintas ( $x_{ikv}$ e $y_{kjv}$ ) para a representação de quantidade de vacinas transportadas.
A separação dessas variáveis é necessária, pois os fluxos de entrada e saída de um centro intermediário podem apresentar distribuições diferentes mesmo que esteja sendo considerado a conservação do fluxo entre etapas. Por exemplo, para um determinado tipo de vacina $v$ em um centro $k$, pode-se ter:

\begin{itemize}
\item Fluxo de entrada:
\[
x_{1kt} = 10, \quad x_{2kt} = 5 \quad \Rightarrow \quad \sum_{i \in I} x_{ikv} = 15
\]

\item Fluxo de saída:
\[
y_{k3t} = 8, \quad y_{k4t} = 7 \quad \Rightarrow \quad \sum_{j \in J} y_{kjv} = 15
\]
\end{itemize}

Observa-se que, embora os fluxos individuais de entrada e saída sejam distintos, o total de vacinas que entra no centro é igual ao total que sai. Dessa forma, a separação entre $x_{ikv}$ e $y_{kjv}$ permite refletir a realidade operacional do sistema logístico.

De forma análoga ao modelo intermunicipal sem intermediação estadual, temos:

\[
\tau_{ik}^{rod}, \quad \tau_{ik}^{areo}, \tau_{kj}^{rod}, \quad \tau_{kj}^{areo} \in \mathbb{R}_{\geq 0}
\]

onde:
\begin{itemize}
    \item $\tau_{ik}^{rod}$ é o tempo de deslocamento do modal rodoviário entre o município de origem $i$ e o centro estadual de destino $k$.
    \item $\tau_{ik}^{areo}$ é o tempo de deslocamento do modal aéreo entre o município de origem $i$ e o centro estadual de destino $k$.
    \item $\tau_{kj}^{rod}$ é o tempo de deslocamento do modal rodoviário entre o centro estadual de origem $k$ e o município de destino $j$.
    \item $\tau_{kj}^{areo}$ é o tempo de deslocamento do modal aéreo entre o centro estadual de origem $k$ e o município de destino $j$.
\end{itemize}

Sejam:
\begin{itemize}
    \item $C_{ik}^{rod}$, o \textbf{Custo Total do transporte rodoviário da primeira etapa};
    \item $C_{ik}^{areo}$, o \textbf{Custo Total do transporte aéreo  da primeira etapa}.
    \item $C_{kj}^{rod}$, o \textbf{Custo Total do transporte rodoviário da segunda etapa};
    \item $C_{kj}^{areo}$, o \textbf{Custo Total do transporte aéreo  da segunda etapa}.
\end{itemize}


Dessa forma, o \textbf{Custo Total de transporte por modal por etapa} é dado por:

\[
C_{ik}^{rod} = \tau_{ik}^{rod} + \delta^{rod}
\]
\[
C_{ik}^{areo} = \alpha \cdot \tau_{ik}^{areo} + \delta^{areo}
\]
\[
C_{kj}^{rod} = \tau_{kj}^{rod} + \delta^{rod}
\]
\[
C_{kj}^{areo} = \alpha \cdot \tau_{kj}^{areo} + \delta^{areo}
\]


Dessa forma, define-se o \textbf{Custo Total de transporte por Etapa} como:

\[
C_{ik} = \min \left( \quad C_{ik}^{rod} \quad , \quad C_{ik}^{areo} \quad \right),  \quad C_{ik} \in \mathbb{R}_{\geq 0}
\]
\[
C_{kj} = \min \left( \quad C_{kj}^{rod} \quad , \quad C_{kj}^{areo} \quad \right),  \quad C_{kj} \in \mathbb{R}_{\geq 0}
\]

Da mesma forma que no modelo anterior, é necessário considerar que a \textbf{capacidade de transporte} varia de acordo com o modal.

Para a \textbf{primeira etapa} (município $i$ $\rightarrow$ estado $k$), sejam:

\[
u_{ik}^{\text{rod}}, \quad u_{ik}^{\text{areo}} \in \mathbb{R}_{\geq 0}
\]

onde:
\begin{itemize}
    \item $u_{ik}^{\text{rod}}$ representa a capacidade máxima agregada de transporte entre o município $i$ e o centro estadual $k$ via modal rodoviário;
    \item $u_{ik}^{\text{areo}}$ representa a capacidade máxima agregada de transporte entre o município $i$ e o centro estadual $k$ via modal aéreo.
\end{itemize}

A capacidade efetiva da primeira etapa é então definida de forma consistente com o custo escolhido:

\begin{equation}
u_{ik} =
\begin{cases}
u_{ik}^{\text{rod}}, & \text{se } C_{ik}^{\text{rod}} \le C_{ik}^{\text{areo}} \\
u_{ik}^{\text{areo}}, & \text{caso contrário}
\end{cases}
\end{equation}

De forma análoga, para a \textbf{segunda etapa} (estado $k$ $\rightarrow$ município $j$), sejam:

\[
u_{kj}^{\text{rod}}, \quad u_{kj}^{\text{areo}} \in \mathbb{R}_{\geq 0}
\]

onde:
\begin{itemize}
    \item $u_{kj}^{\text{rod}}$ representa a capacidade máxima de transporte entre o centro estadual $k$ e o município $j$ via modal rodoviário;
    \item $u_{kj}^{\text{areo}}$ representa a capacidade máxima de transporte entre o centro estadual $k$ e o município $j$ via modal aéreo.
\end{itemize}

A capacidade efetiva da segunda etapa é definida por:

\begin{equation}
u_{kj} =
\begin{cases}
u_{kj}^{\text{rod}}, & \text{se } C_{kj}^{\text{rod}} \le C_{kj}^{\text{areo}} \\
u_{kj}^{\text{areo}}, & \text{caso contrário}
\end{cases}
\end{equation}


Quanto ao tempo de distribuição de vacinas, o tempo total de entrega de uma dose corresponderia, conceitualmente, à soma dos tempos das duas etapas do transporte: do município de origem ao centro estadual e do centro estadual ao município de destino.

Entretanto, a modelagem adotada não permite o rastreamento explícito de cada unidade de vacina ao longo de todo o percurso, já que as variáveis de decisão representam fluxos agregados e isso impede de associar os trechos das duas etapas a uma mesma unidade transportada. Dessa forma, não é possível determinar diretamente o tempo exato de distribuição de uma dose específica. 

Em vez disso, opta-se por calcular, em um pós-processamento da solução ótima, um \textbf{limite superior} para o tempo total de distribuição, a partir da composição de arcos com fluxo positivo.

Define-se, portanto, o \textbf{Limite Superior do Tempo de Distribuição} como:

\[
L = \max \left\{ \; C_{ik} + C_{kj} 
\;\middle|\;
x_{ikv} > 0,\; y_{kjv} > 0
\; \right\}
\]

onde:
\begin{itemize}
    \item $L$ representa um limite superior para o tempo total de distribuição;
    \item $C_{ik}$ é o tempo de transporte entre o município de origem $i$ e o centro estadual $k$;
    \item $C_{kj}$ é o tempo de transporte entre o centro estadual $k$ e o município de destino $j$;
    \item $x_{ikv}$ e $y_{kjv}$ são as variáveis de decisão associadas aos fluxos nas duas etapas do sistema.
\end{itemize}

Essa formulação não garante que os dois trechos considerados pertençam ao mesmo caminho físico percorrido por uma única dose, uma vez que não há rastreamento individual das unidades. Assim, o valor de $L$ deve ser interpretado como um \textbf{limite superior}, isto é, uma estimativa conservadora do pior tempo possível compatível com os fluxos positivos da solução ótima.

Ainda assim, esse limite superior fornece uma aproximação útil do tempo total de distribuição e pode ser utilizado como uma \textbf{referência de desempenho} para avaliar a eficiência temporal do sistema logístico.


\subsection{Função Objetivo}

Dessa forma, a \textbf{função objetivo} fica definida:

\begin{equation}
    \min z = \sum_{v} \sum_{i} \sum_{k} C_{ik} \cdot x_{ikv} \quad + \quad \sum_{v} \sum_{k} \sum_{j} C_{kj} \cdot y_{kjv}
    \label{eq:intermunicipal_estendido}
\end{equation}

ou, com a penalidade:

\begin{equation}
    \min z = \sum_{v} \sum_{i} \sum_{k} C_{ik} \cdot x_{ikv} \cdot \rho_{iv} \quad + \quad \sum_{v} \sum_{k} \sum_{j} C_{kj} \cdot y_{kjv}
    \label{eq:intermunicipal_estendido}
\end{equation}

onde a primeira parcela representa o custo da primeira etapa do fluxo e a segunda parcela representa o custo da segunda etapa.

A penalidade é aplicada na origem, pois o risco de expiração está associado ao tempo restante no município de origem.

\subsection{Restrições}

Para garantir a coerência geográfica e administrativa do fluxo logístico, é necessário impor restrições adicionais que assegurem que a intermediação estadual ocorra exclusivamente por meio do centro estadual $k$ pertencente ao estado de destino do município $j$ atendido.

Para isso, sejam $\text{estado}(k)$ e $\text{estado}(j)$ funções que retornam o estado associado ao centro estadual $k$ e ao município $j$, respectivamente.

A \textbf{Restrição de Consistência Estadual} é dada por:

\[
y_{kjv} = 0, \quad \forall (k, j) \in K \times J \mid \text{estado}(k) \neq \text{estado}(j), \forall v \in V
\]

Essa restrição bloqueia a saída de vacinas de centros estaduais {k} cujo estado seja diferente do estado do município {j}, forçando-se a escolher um centro diferente até que seja escolhido um centro {k} do mesmo estado que $j$.

A \textbf{restrição de oferta} garante que a quantidade total de vacinas enviadas por um município de origem não exceda o estoque disponível daquele município para cada tipo de vacina.

\begin{equation*}
\sum_{k} x_{ikv} \le o_{ik}, \quad \forall i \in I,\; v \in V
\end{equation*}

A \textbf{restrição de demanda} assegura que toda a demanda por vacinas em cada município de destino seja completamente atendida.

\begin{equation*}
\sum_{k} y_{kjv} = d_{jv}, \quad \forall j \in J, \; v \in V
\end{equation*}

A \textbf{restrição de conservação de fluxo} garante que, para cada centro intermediário, a quantidade total de vacinas recebidas seja igual à quantidade total redistribuída, não havendo acúmulo ou perda de vacinas no centro.

\begin{equation*}
\sum_{i} x_{ikv} = \sum_{j} y_{kjv}, \quad \forall k \in K,\; v \in V
\end{equation*}

A \textbf{primeira restrição de capacidade agregada} é relativa a capacidade de transporte da primeira etapa do fluxo e, nesta restrição, a quantidade de vacinas transportadas não pode ultrapassar esse limite. Logo, temos:
\begin{equation*}
\sum_{v} x_{ikv} \le u_{ik}, \quad \forall i \in I, \; k \in K
\end{equation*}
onde:
\begin{itemize}
    \item $u_{ik}$ representa a capacidade de transporte do veículo na primeira etapa do fluxo.
\end{itemize}

A \textbf{segunda restrição de capacidade agregada} é relativa a capacidade de transporte da segunda etapa do fluxo e, nesta restrição, a quantidade de vacinas transportadas não pode ultrapassar esse limite. Logo, temos:
\begin{equation*}
\sum_{v} y_{kjv} \le u_{kj}, \quad \forall k \in K, \; j \in J
\end{equation*}
onde:
\begin{itemize}
    \item $u_{kj}$ representa a capacidade de transporte do veículo na segunda etapa do fluxo.
\end{itemize}

A \textbf{restrição de viabilidade temporal} é relativa à ideia de remover pares onde o tempo de transporte é maior que o tempo restante até o vencimento. Assim, temos:

\begin{equation*}
    x_{ikv} = 0, \quad \forall i \in I, \; k \in K, \; v \in V \text{ tais que } C_{ik} > e_{iv}
\end{equation*}

A restrição de viabilidade temporal é aplicada apenas na primeira etapa do fluxo. Caso o tempo de transporte entre o município de origem e o centro estadual exceda o tempo restante até a expiração da vacina, o envio é proibido.

Não é necessário impor restrições adicionais sobre os fluxos de saída dos centros estaduais, pois a restrição de conservação de fluxo garante que, na ausência de entrada viável, não haverá saída correspondente, assegurando a consistência do modelo.


A \textbf{restrição de não-negatividade} assegura que as quantidades de vacinas transportadas não podem ser negativas e devem ser números inteiros.
\begin{equation*}
x_{ikv}, y_{kjv} \in \mathbb{Z}_{\geq 0}, \quad \forall i \in I, k \in K, j \in J, v \in V
\end{equation*}

Em suma, o PPLI ficou definido assim:

\begin{equation*}
    \min z = \sum_{v} \sum_{i} \sum_{k} C_{ik} \cdot x_{ikv} \quad + \quad \sum_{v} \sum_{k} \sum_{j} C_{kj} \cdot y_{kjv}
\end{equation*}

sujeito a:

\begin{equation*}
y_{kjv} = 0, \quad \forall (k, j) \in K \times J \mid \text{estado}(k) \neq \text{estado}(j), \forall v \in V
\end{equation*}

\begin{equation*}
\sum_{k} x_{ikv} \le o_{ik}, \quad \forall i \in I,\; v \in V
\end{equation*}

\begin{equation*}
\sum_{k} y_{kjv} = d_{jv}, \quad \forall j \in J, \; v \in V
\end{equation*}

\begin{equation*}
\sum_{i} x_{ikv} = \sum_{j} y_{kjv}, \quad \forall k \in K,\; v \in V
\end{equation*}

\begin{equation*}
\sum_{v} x_{ikv} \le u_{ik}, \quad \forall i \in I, \; k \in K
\end{equation*}

\begin{equation*}
\sum_{v} y_{kjv} \le u_{kj}, \quad \forall k \in K, \; j \in J
\end{equation*}

\begin{equation*}
    x_{ikv} = 0, \quad \forall i \in I, \; k \in K, \; v \in V \text{ tais que } C_{ik} > e_{iv}
\end{equation*}

\begin{equation*}
x_{ikv}, y_{kjv} \in \mathbb{Z}_{\geq 0}, \quad \forall i \in I, k \in K, j \in J, v \in V
\end{equation*}

\bigskip

\section{Modelagem do transporte intramunicipal de vacinas}

A distribuição intramunicipal de vacinas corresponde à etapa final da cadeia logística, na qual as doses são transportadas da secretaria municipal de saúde para os pontos de aplicação, como postos de saúde e hospitais, ou seja, existe apenas um ponto de origem das vacinas.

Para o caso intramunicipal, somente o modal rodoviário é considerado.

\subsection{Variável de Decisão e Parâmetros}

Seja:

\[
z_{jpv} \in \mathbb{Z}_{\geq 0}
\]

onde:
\begin{itemize}
    \item $z_{jpv}$ representa a quantidade de vacinas do tipo $v$ transportadas do centro do município $j$ para o posto de destino $p$;
    \item $j \in J$ representa os municípios;
    \item $p \in P_j$ representa os postos de saúde pertencentes ao município $j$;
    \item $v \in V$ representa os tipos de vacina.
\end{itemize}

Como se trata de unidades discretas (doses), temos:

\begin{equation*}
z_{jpv} \in \mathbb{Z}_{\geq 0}, \quad \forall j \in J, p \in P_j, v \in V
\end{equation*}

De forma análoga ao modelo intermunicipal, temos:

\[
\tau_{jp}^{rod} \in \mathbb{R}_{\geq 0}
\]
\[
\delta^{rod} \in \mathbb{R}_{\geq 0}
\]

onde:
\begin{itemize}
    \item $\tau_{jp}^{rod}$ é o tempo de deslocamento do modal rodoviário entre o centro de distribuição municipal $j$ e o posto de destino $p$.
    \item $\delta^{rod}$ representa o tempo de processamento logístico do modal rodoviário.
\end{itemize}

Seja:

\begin{equation*}
C_{jp} \in \mathbb{R}_{\geq 0}
\end{equation*}

o \textbf{Custo Total do transporte} entre os pontos $j$ e $p$.

Dessa forma, o \textbf{Custo Total de transporte} é dado por:

\[
C_{jp} = \tau_{jp}^{rod} + \delta^{rod}
\]

No contexto intramunicipal, considera-se que o transporte é realizado exclusivamente por meio do modal rodoviário. Dessa forma, a capacidade de transporte não depende da escolha entre diferentes modais.

Seja:

\[
q_{jp} \in \mathbb{R}_{\geq 0}
\]

onde:
\begin{itemize}
    \item $q_{jp}$ representa a capacidade máxima de transporte entre o centro de distribuição do município $j$ e o posto de destino $p$.
\end{itemize}

Seguindo um raciocínio análogo ao do modelo intermunicipal sem intermediação estadual, o \textbf{tempo de distribuição} das vacinas é dado por:

\[
T = \max \left( \quad C_{jp} \quad \mid \quad z_{jpv} > 0 \quad \right)
\]

\subsection{Função Objetivo}

A \textbf{função objetivo} é definida como:

\begin{equation}
\min z = \sum_{v} \sum_{j} \sum_{p \in P_j} C_{jp} \cdot z_{jpv}
\label{intramunicipal_sem_penalidade}
\end{equation}

O objetivo é minimizar o tempo total de distribuição das vacinas dentro do município.

Para incorporar o risco de expiração das vacinas, introduz-se um fator de penalização associado ao tempo restante de validade no município. Seja:

\[
\rho_{jv} = \frac{1}{e_{jv}}
\]
onde:
\begin{itemize}
\item  $e_{jv}$ representa o tempo restante até a expiração da vacina do tipo $v$ disponível no município $j$.
\item $\rho_{jv}$ é penalidade da vacina do tipo $v$ no município $j$
\end{itemize}

Dessa forma, a função objetivo passa a ser:
\begin{equation}
\min z = \sum_{v} \sum_{j} \sum_{p \in P_j} C_{jp} \cdot z_{jpv} \cdot \rho_{jv}
\label{intramunicipal_com_penalidade}
\end{equation}

Essa formulação prioriza a distribuição de vacinas com menor tempo restante de validade associado ao estoque no nível intramunicipal.

\subsection{Restrições}

A \textbf{restrição de oferta} diz respeito a quantidade total enviada por cada ponto de origem não poder exceder sua disponibilidade:

\begin{equation*}
\sum_{p \in P_j} z_{jpv} \le s_{jv}, \quad \forall j \in J, \; v \in V
\end{equation*}

onde:
\begin{itemize}
    \item $s_{jv}$ representa a quantidade disponível da vacina $t$ no ponto de origem $j$.
\end{itemize}

A \textbf{restrição de demanda} estabelece que cada ponto de destino deve receber exatamente sua demanda:

\begin{equation*}
\sum_{j} z_{jpv} = r_{pv}, \quad \forall p \in P_j, \; v \in V
\end{equation*}

onde:
\begin{itemize}
    \item $r_{pv}$ representa a demanda da vacina $t$ no ponto de destino $p$.
\end{itemize}

A \textbf{restrição de capacidade agregada} considera que a capacidade de transporte entre os pontos é limitada:

\begin{equation*}
\sum_{v} z_{jpv} \le q_{jp}, \quad \forall j \in J, p \in P_j
\end{equation*}

onde:
\begin{itemize}
    \item $q_{jp}$ representa a capacidade máxima de transporte entre $j$ e $p$.
\end{itemize}

A \textbf{restrição de viabilidade} temporal é definida para evitar perdas por expiração:

\begin{equation*}
z_{jpv} = 0, \quad \forall j \in J, \; p \in P_j, \; v \in V \text{ tais que } C_{jp} > e_{jv}
\end{equation*}

onde:
\begin{itemize}
    \item $e_{jv}$ representa o tempo restante até a expiração da vacina $v$ no ponto de origem $j$.
\end{itemize}

O modelo completo é dado por:

\begin{equation*}
\min z = \sum_{v} \sum_{j} \sum_{b} C_{jp} \cdot z_{jpv}
\end{equation*}
sujeito a:
\begin{equation*}
\sum_{b} z_{jpv} \le s_{jv}, \quad \forall j \in J, \; v \in V
\end{equation*}
\begin{equation*}
\sum_{j} z_{jpv} = r_{pv}, \quad \forall p \in P_j, \; v \in V
\end{equation*}
\begin{equation*}
\sum_{v} z_{jpv} \le q_{jp}, \quad \forall j \in J, p \in P_j
\end{equation*}
\begin{equation*}
z_{jpv} = 0, \quad \forall j \in J, \; p \in P_j, \; v \in V \text{ tais que } C_{jp} > e_{jv}
\end{equation*}
\begin{equation*}
z_{jpv} \in \mathbb{Z}_{\geq 0}, \quad \forall j \in J, p \in P_j, v \in V
\end{equation*}


\bigskip

\section{Modelo Unificado Sem Intermediação Estadual}

Este modelo considera transportes direto entre municípios sem o intermédio de órgãos estaduais.

A ideia de unificação representa a integração dos modelos intermunicipal e intramunicipal em uma única formulação matemática. O objetivo é representar, de forma unificada, todo o fluxo logístico das vacinas, desde a redistribuição entre municípios até a entrega final nos postos de vacinação. A principal motivação dessa abordagem é permitir a otimização global do sistema.

\subsection{Variáveis de Decisão e Parâmetros}

O modelo unificado considera dois níveis de decisão:

\begin{itemize}
    \item Transporte intermunicipal:
    \[
    x_{ijv}
    \]
    onde $x_{ijv}$ representa a quantidade de vacinas do tipo $v$ transportadas do município $i$ para o município $j$.
    
    \item Transporte intramunicipal:
    \[
    z_{jpv}
    \]
    onde $z_{jpv}$ representa a quantidade de vacinas do tipo $v$ transportadas do ponto de origem $j$ para o posto de destino $p$ dentro de um município.
\end{itemize}

Neste modelo, a distribuição de vacinas pode ocorrer em duas etapas distintas: uma etapa intermunicipal e outra intramunicipal.

Embora o tempo total de entrega de uma dose de vacina, do ponto de origem até o destino final, seja conceitualmente dado pela soma dos tempos das etapas intermunicipal e intramunicipal, o modelo proposto não permite o rastreamento explícito do percurso completo de cada unidade transportada, já que as variáveis de decisão representam fluxos agregados.

Dessa forma, não é possível determinar diretamente, durante a otimização, o tempo exato de distribuição de uma dose específica. Em vez disso, opta-se por calcular, em um pós-processamento da solução ótima, um \textbf{limite superior} para o tempo total de distribuição, a partir da composição de arcos com fluxo positivo.

Define-se, portanto, o \textbf{limite superior do tempo de distribuição} como:

\[
L = \max \left\{ \; C_{ij} + C_{jp} 
\;\middle|\;
x_{ijv} > 0,\; z_{jpv} > 0
\; \right\}
\]

onde:
\begin{itemize}
    \item $L$ representa um limite superior para o tempo total de distribuição das vacinas;
\end{itemize}

Ressalta-se que essa formulação não garante que os dois trechos considerados pertençam ao mesmo caminho físico percorrido por uma única dose, uma vez que não há rastreamento individual das unidades. Assim, o valor de $L$ deve ser interpretado como um \textbf{limite superior}, isto é, uma estimativa conservadora do pior tempo possível compatível com os fluxos positivos da solução ótima.

Ainda assim, esse limite superior fornece uma aproximação útil do tempo total de distribuição e pode ser utilizado como uma \textbf{referência de desempenho} para avaliar a eficiência temporal do sistema logístico.

Com a unificação, torna-se possível estabelecer uma relação entre os tempos de expiração $e_{iv}$ e $e_{jv}$. Intuitivamente, $e_{jv}$ pode ser interpretado como o tempo remanescente de expiração das vacinas do tipo $v$ disponíveis no município $j$ após a etapa de redistribuição intermunicipal.

No entanto, devido à estrutura do modelo, não é possível rastrear explicitamente a origem de cada unidade de vacina que chega ao município $j$ com as variáveis de decisão. Ou seja, não se sabe exatamente quanto cada município $i$ contribuiu para o estoque final de $j$. Como consequência, não é possível calcular de forma exata o tempo de expiração remanescente com base nas variáveis de decisão do modelo.

Diante dessa limitação, define-se $e_{jv}$ como um parâmetro estimado \textit{a priori}, isto é, calculado antes da resolução do modelo. A ideia adotada é considerar que o conjunto de vacinas que chega ao município $j$ representa uma mistura das vacinas disponíveis nos municípios de origem. Assim, assume-se que essa mistura reflete, em termos agregados, a distribuição das ofertas iniciais $o_{iv}$ na rede.
Além disso, deve-se considerar que as vacinas perdem tempo de validade durante o transporte. Dessa forma, uma vacina que sai de $i$ com validade $e_{iv}$ chega a $j$ com validade reduzida, tem seu tempo de expiração reduzido por:

\[
e_{iv} - C_{ij}
\]

Com base nessas hipóteses, define-se $e_{jv}$ como uma média ponderada dos tempos remanescentes de expiração das vacinas provenientes de cada município de origem, utilizando como pesos as ofertas iniciais $o_{iv}$. Assim, tem-se:

\begin{equation*}
e_{jv} =
\frac{
\sum_{i \in I} o_{iv} \cdot \left( e_{iv} - C_{ij} \right)
}{
\sum_{i \in I} o_{iv}
},
\quad \forall j \in J,\; v \in V
\end{equation*}

\subsection{Função Objetivo}

A \textbf{função objetivo} unificada é dada por:

\begin{equation}
\min z = 
\sum_{v} \sum_{i} \sum_{j} C_{ij} \cdot x_{ijv}
\;+\;
\sum_{v} \sum_{j} \sum_{p \in P_j} C_{jp} \cdot z_{jpv}
\label{unificado_sem_penalidade}
\end{equation}

A primeira parcela representa o custo do transporte intermunicipal, enquanto a segunda representa o custo da distribuição intramunicipal.

Para incorporar o risco de expiração, pode-se adicionar os fatores de penalidade e a função objetivo passa a ser:

\begin{equation}
\min z = 
\sum_{v} \sum_{i} \sum_{j} C_{ij} \cdot x_{ijv} \cdot \rho_{iv}
\;+\;
\sum_{v} \sum_{j} \sum_{p \in P_j} C_{jp} \cdot z_{jpv} \cdot \rho_{jv}
\label{unificado_com_penalidade}
\end{equation}

Essa formulação prioriza o transporte de vacinas com menor tempo restante até a expiração em ambos os níveis da cadeia logística.

\subsection{Restrições}

As \textbf{Restrições Intermunicipais} a serem reaproveitadas são:

\begin{equation*}
\sum_{j} x_{ijv} \le o_{iv}, \quad \forall i \in I, \; v \in V
\end{equation*}

\begin{equation*}
\sum_{v} x_{ijv} \le u_{ij}, \quad \forall i \in I, j \in J
\end{equation*}

\begin{equation*}
x_{ijv} \in \mathbb{Z}_{\geq 0}
\end{equation*}

A restrição de demanda no nível intermunicipal foi removida no modelo unificado, uma vez que a demanda final por vacinas é imposta explicitamente no nível intramunicipal, por meio dos postos de vacinação, ou seja, a demanda de um município é dada pela soma da demanda dos seus postos de saúde.

A manutenção simultânea da restrição de demanda no nível intermunicipal e no nível intramunicipal introduziria redundância e potencial inconsistência no modelo, podendo torná-lo inviável caso os valores agregados não coincidam exatamente.

As \textbf{Restrições Intramunicipais} a serem reaproveitadas são:

\begin{equation*}
\sum_{j} z_{jpv} = r_{pv}, \quad \forall p \in P_j, \; v \in V
\end{equation*}

\begin{equation*}
\sum_{v} z_{jpv} \le q_{jp}, \quad \forall j \in J, p \in P_j
\end{equation*}

\begin{equation*}
z_{jpv} \in \mathbb{Z}_{\geq 0}
\end{equation*}

A restrição de oferta no nível intramunicipal foi removida devido à reformulação da estrutura de acoplamento entre os níveis intermunicipal e intramunicipal.

No modelo unificado, a quantidade de vacinas disponível para distribuição dentro de cada município é determinada endogenamente pela soma do estoque inicial local com o fluxo recebido de outros municípios. Dessa forma, a imposição de uma restrição adicional limitando a saída com base em um parâmetro fixo de oferta tornaria o modelo inconsistente, pois poderia restringir indevidamente o fluxo de vacinas disponível para atendimento da demanda final.

As \textbf{restrições de viabilidade temporal} são aplicadas separadamente em cada nível da cadeia logística. Essa abordagem decorre do fato de que o modelo não rastreia explicitamente o percurso completo de cada unidade de vacina desde o município de origem até o ponto de vacinação final.

Dessa forma, não é possível impor diretamente uma restrição baseada na soma dos tempos das etapas intermunicipal e intramunicipal para uma mesma unidade transportada e que seja da forma:

\[
C_{ij} + C_{jp} \le e_{iv}
\]

Assim, adota-se uma aproximação em que a viabilidade temporal é garantida localmente em cada etapa do processo. Isso assegura que as vacinas não permaneçam em trânsito por um tempo superior ao seu prazo de validade em nenhuma das etapas, mantendo a coerência com a estrutura do modelo e sua tratabilidade computacional. Dessa forma, temos:

\begin{equation*}
x_{ijv} = 0, \quad \forall i \in I, \; j \in J, \; v \in V \text{ tais que } C_{ij} > e_{iv}
\end{equation*}

\begin{equation*}
z_{jpv} = 0, \quad \forall j \in J, \; p \in P_j, \; v \in V \text{ tais que } C_{jp} > e_{jv}
\end{equation*}

A integração entre os níveis intermunicipal e intramunicipal é consolidada por meio da \textbf{Restrição de Conservação e Disponibilidade de Estoque} que garante que o sistema respeite a realidade física do estoque em cada localidade.

Esta restrição descreve o que acontece com o estoque de vacinas de cada município que pode receber vacinas de outras localidades ou fornece-las para outros municípios ou postos. Logo, a lógica matemática estabelece que a soma de tudo o que sai do município (seja para o consumo interno em seus postos ou para exportação a outras localidades) não pode ser superior à soma de tudo o que o município dispõe (seu estoque inicial somado ao que ele recebeu de fora durante o processo). 

Dessa forma, a restrição é expressa como:

\begin{equation*}
\sum_{p \in P_j} z_{jpv} + \sum_{j' \neq j} x_{jj'v} \le o_{jv} + \sum_{i \neq j} x_{ijv}, \quad \forall j \in J, \; v \in V
\end{equation*}

Esta formulação elimina a necessidade de variáveis intermediárias de estoque final, pois o balanço é verificado globalmente. Além de prevenir a criação artificial de doses no modelo, o uso de uma desigualdade ($\le$) permite que o sistema opte por manter parte das vacinas armazenadas se não houver demanda imediata, refletindo a flexibilidade necessária em operações logísticas reais.


Em suma, o PPLI fica definido assim:

\begin{equation*}
\min z = 
\sum_{v} \sum_{i} \sum_{j} C_{ij} \cdot x_{ijv}
\;+\;
\sum_{v} \sum_{j} \sum_{p \in P_j} C_{jp} \cdot z_{jpv}
\end{equation*}

sujeito a:

\begin{equation*}
\sum_{j} x_{ijv} \le o_{iv}, \quad \forall i \in I, \; v \in V
\end{equation*}

\begin{equation*}
\sum_{v} x_{ijv} \le u_{ij}, \quad \forall i \in I, j \in J
\end{equation*}

\begin{equation*}
\sum_{j} z_{jpv} = r_{pv}, \quad \forall p \in P_j, \; v \in V
\end{equation*}

\begin{equation*}
\sum_{v} z_{jpv} \le q_{jp}, \quad \forall j \in J, p \in P_j
\end{equation*}

\begin{equation*}
\sum_{p \in P_j} z_{jpv} + \sum_{j' \neq j} x_{jj'v} \le o_{jv} + \sum_{i \neq j} x_{ijv}, \quad \forall j \in J, \; v \in V
\end{equation*}

\begin{equation*}
x_{ijv} = 0, \quad \forall i \in I, \; j \in J, \; v \in V \text{ tais que } C_{ij} > e_{iv}
\end{equation*}

\begin{equation*}
z_{jpv} = 0, \quad \forall j \in J, \; p \in P_j, \; v \in V \text{ tais que } C_{jp} > e_{jv}
\end{equation*}

\begin{equation*}
x_{ijv}, z_{jpv} \in \mathbb{Z}_{\geq 0}
\end{equation*}

\section{Modelo Unificado Com Intermediação Estadual}

Para representar a intermediação estadual, incorpora-se um nível adicional na modelagem intermunicipal.

\subsection{Variáveis de Decisão e Parâmetros}

\begin{itemize}
    \item $x_{ikv}$: fluxo do município $i$ para o centro estadual $k$;
    \item $y_{kjv}$: fluxo do centro estadual $k$ para o município $j$;
    \item $z_{jpv}$: fluxo intramunicipal.
\end{itemize}

Da mesma forma que modelo anterior, do ponto de vista operacional, o tempo total de entrega de uma dose de vacina corresponderia à soma dos tempos dessas três etapas. No entanto, de forma análoga ao modelo sem intermediação estadual, a estrutura matemática adotada não permite o rastreamento explícito do fluxo completo de cada unidade ao longo de toda a cadeia logística, pois não existe uma variável que represente simultaneamente o percurso completo de uma vacina desde a origem até o destino final. Como consequência, não é possível garantir que os fluxos modelados em cada etapa correspondam à mesma unidade transportada, inviabilizando a soma direta dos tempos por trajeto completo.

Dessa forma, opta-se por uma abordagem consistente com a modelagem proposta, na qual o tempo total de distribuição é definido como a soma dos piores tempos de cada etapa.


Assim, o \textbf{tempo total de distribuição} é dado por:

\[
T = \max \left\{ \; C_{ik} \;\middle|\; x_{ikv} > 0 \; \right\} + 
\max \left\{ \; C_{kj} \;\middle|\; y_{kjv} > 0 \; \right\} + 
\max \left\{ \; C_{jp} \;\middle|\; z_{jpv} > 0 \; \right\}
\]

onde:
\begin{itemize}
    \item $T$ representa o tempo total de distribuição das vacinas;
    \item $C_{ik}$ é o custo (tempo) do transporte entre o município de origem e o centro estadual;
    \item $C_{kj}$ é o custo (tempo) do transporte entre o centro estadual e o município de destino;
    \item $C_{jp}$ é o custo (tempo) do transporte intramunicipal.
\end{itemize}

Além disso, como no modelo anterior, torna-se possível estabelecer uma relação entre os tempos de expiração $e_{iv}$ e $e_{jv}$. Diante da impossibilidade de rastreio explícito de cada unidade de vacina, $e_{jv}$ também será definido como um parâmetro estimado \textit{a priori}.

Como no modelo com intermediação estadual a degradação temporal das vacinas ocorre em duas etapas, será necessário dividir o cálculo de $e_{jv}$ entre essas duas etapas.

Na primeira etapa, por um raciocínio análogo ao do modelo anterior, define-se o tempo médio de expiração no centro estadual $k$ como uma média ponderada dos tempos remanescentes de expiração das vacinas provenientes de cada município de origem, utilizando como pesos as ofertas iniciais $o_{iv}$. Além disso, considera-se a perda de validade ao longo do transporte dada por ($e_{iv} - C_{ik}$).

Dessa forma, define-se o tempo médio de expiração no centro estadual $k$ como:

\begin{equation*}
e_{kv} =
\frac{
\sum_{i \in I} o_{iv} \cdot \left( e_{iv} - C_{ik} \right)
}{
\sum_{i \in I} o_{iv}
},
\quad \forall k \in K,\; v \in V
\end{equation*}

Na segunda etapa, o problema se torna mais delicado. Diferentemente da primeira etapa, na qual é possível relacionar diretamente os municípios de origem $i$ com o ponto intermediário $k$, aqui não existe mais uma ligação explícita entre os municípios de origem $i$ e o município de destino final $j$.

Intuitivamente, as vacinas que chegam ao município $j$ já passaram por um processo de agregação no centro estadual $k$. Nesse processo, vacinas provenientes de diferentes origens são misturadas, perdendo-se completamente a informação sobre suas origens individuais. Como consequência, não é mais possível utilizar diretamente as ofertas iniciais $o_{iv}$ como pesos, nem reconstruir com precisão a contribuição de cada origem no estoque final de $j$.

Para contornar essa limitação, adota-se uma aproximação baseada na ideia de que os centros estaduais mais próximos (com menor tempo de transporte $C_{kj}$) do município de destino $j$ recebem um peso matematicamente maior ($\frac{1}{C_{kj}}$) no cálculo da validade média que chega ao município. Além disso, considera-se novamente a perda de validade ao longo do transporte até o município $j$ dada por ($e_{kv} - C_{kj}$).

Dessa forma, define-se o tempo médio de expiração no município $j$ como

\begin{equation*}
e_{jv} =
\frac{
\sum_{k \in K} \frac{1}{C_{kj}} \cdot \left( e_{kv} - C_{kj} \right)
}{
\sum_{k \in K} \frac{1}{C_{kj}}
},
\quad \forall j \in J,\; v \in V
\end{equation*}

Além disso, assume-se que $C_{kj} > 0$, uma vez que sempre existe um custo temporal associado ao deslocamento entre o centro estadual e o município de destino.

\subsection{Função Objetivo}

\begin{equation}
\min z =
\sum_{v} \sum_{i} \sum_{k} C_{ik} \cdot x_{ikv}
+
\sum_{v} \sum_{k} \sum_{j} C_{kj} \cdot y_{kjv}
+
\sum_{v} \sum_{j} \sum_{p \in P_j} C_{jp} \cdot z_{jpv}
\label{unificado_estendido}
\end{equation}

ou, com as penalidades:

\begin{equation}
\min z =
\sum_{v} \sum_{i} \sum_{k} C_{ik} \cdot x_{ikv} \cdot \rho_{iv}
+
\sum_{v} \sum_{k} \sum_{j} C_{kj} \cdot y_{kjv}
+
\sum_{v} \sum_{j} \sum_{p \in P_j} C_{jp} \cdot z_{jpv} \cdot \rho_{jv}
\label{unificado_estendido}
\end{equation}

\subsection{Restrições}

Seguindo um raciocínio análogo do modelo anterior, as \textbf{restrições de viabilidade temporal} são aplicadas separadamente em cada etapa da cadeia logística (municipal, estadual e intramunicipal). Essa abordagem decorre do fato de que o modelo não rastreia explicitamente o percurso completo de cada unidade de vacina ao longo das três etapas.

A \textbf{Restrição de Conservação e Disponibilidade de Estoque} também deve ser implementada neste modelo seguindo um raciocínio análogo do modelo anterior e fazendo-se as adaptações necessárias considerando o intermédio dos centros estaduais.

Dessa forma, as restrições são:

\[
y_{kjv} = 0, \quad \forall (k, j) \in K \times J \mid \text{estado}(k) \neq \text{estado}(j), \forall v \in V
\]

\begin{equation*}
\sum_{k} x_{ikv} \le o_{iv}, \quad \forall i \in I, \; v \in V
\end{equation*}

\begin{equation*}
\sum_{i} x_{ikv} = \sum_{j} y_{kjv}, \quad \forall k \in K, \; v \in V
\end{equation*}

\begin{equation*}
\sum_{j} z_{jpv} = r_{pv}, \quad \forall p \in P_j, \; v \in V
\end{equation*}

\begin{equation*}
\sum_{v} x_{ikv} \le u_{ik}, \quad \forall i \in I, k \in K
\end{equation*}

\begin{equation*}
\sum_{v} y_{kjv} \le u_{kj}, \quad \forall k \in K, j \in J
\end{equation*}

\begin{equation*}
\sum_{v} z_{jpv} \le q_{jp}, \quad \forall j \in J, p \in P_j
\end{equation*}

\begin{equation*}
y_{kjv} = 0, \quad \forall (k, j) \in K \times J \mid \text{estado}(k) \neq \text{estado}(j), \forall v \in V
\end{equation*}

\begin{equation*}
\sum_{p \in P_j} z_{jpv} + \sum_{k \in K} x_{jkv} \le o_{jv} + \sum_{k \in K} y_{kjv}, \quad \forall j \in J, \; v \in V
\end{equation*}

\begin{equation*}
    x_{ikv} = 0, \quad \forall i \in I, \; k \in K, \; v \in V \text{ tais que } C_{ik} > e_{iv}
\end{equation*}

\begin{equation*}
    z_{jpv} = 0, \quad \forall j \in J, \; p \in P_j, \; v \in V \text{ tais que } C_{jp} > e_{jv}
\end{equation*}

\begin{equation*}
x_{ikv}, y_{kjv}, z_{jpv} \in \mathbb{Z}_{\geq 0}
\end{equation*}

\bigskip